\section{Idee Vincent}

\subsection{Grundkonzept}

Die Idee des Low-Fidelity-Prototyps \textit{Wien in Zahlen} ist ein personalisierter Open-Data-Feed für die Stadt Wien. Die App soll nicht wie ein klassisches Datenportal funktionieren, in dem viele Datensätze gleichzeitig angezeigt werden. Stattdessen erhalten Nutzer:innen regelmäßig kurze und verständlich aufbereitete Datenbeiträge in Form von Fact Sheets. Dadurch verbindet die App Elemente einer News-App mit den Möglichkeiten eines Open-Data-Portals.

\subsection{Personalisierter Feed}

Zu Beginn der Nutzung wählen die Nutzer:innen Themen aus, die sie besonders interessieren, zum Beispiel Umwelt, Wohnen, Mobilität oder Energie. Auf Basis dieser Auswahl wird ein persönlicher Feed erstellt. In diesem Feed erscheinen kurze Karten mit aktuellen oder interessanten Datenbeiträgen zur Stadt Wien. Jede Karte gibt einen ersten Überblick und führt bei Interesse zu einem ausführlicheren Fact Sheet. So können Nutzer:innen selbst entscheiden, welche Inhalte sie genauer ansehen möchten.

\subsection{Aufbau der Fact Sheets}

Ein Fact Sheet enthält eine kurze Erklärung des Themas, eine einfache Datenvisualisierung, die verwendete Datenquelle, einen Vergleich sowie den Abschnitt \textit{Was bedeutet das?}. Dieser Abschnitt ist besonders wichtig, weil er die zentrale Aussage der Daten in Alltagssprache erklärt. Die Nutzer:innen sollen nicht nur Zahlen sehen, sondern auch verstehen, welche Bedeutung diese Zahlen für Wien und den Alltag in der Stadt haben.

\subsection{Unterstützung beim Verstehen von Visualisierungen}

Ergänzend gibt es bei jeder Visualisierung ein \textit{Grafik verstehen}-Element mit einem Fragezeichen. Wenn dieses angeklickt wird, öffnet sich eine kurze Erklärung zur jeweiligen Grafik, zum Beispiel wie ein Balkendiagramm gelesen wird. Dadurch unterstützt die App auch Menschen, die wenig Erfahrung mit Datenvisualisierungen oder Statistiken haben.

\subsection{Voting und Frageportal}

Zusätzlich enthält \textit{Wien in Zahlen} ein Voting-System. Damit können Nutzer:innen abstimmen, welche Themen oder Datenfragen in Zukunft aufbereitet werden sollen. Ein Frageportal ermöglicht es außerdem, eigene Fragen zur Stadt Wien einzureichen. Dadurch wird die App nicht nur zu einem Informationsangebot, sondern auch zu einem Mitmach-Format.

\subsection{Ziel der Umsetzung}

Die Umsetzung ist sinnvoll, weil sie Open Data niedrigschwellig, personalisiert und alltagsnah vermittelt. Komplexe Daten werden in kleine, verständliche Einheiten übersetzt, wodurch mehr Menschen einen Zugang zu städtischen Informationen erhalten und sich besser mit Entwicklungen in Wien auseinandersetzen können.